\section{Diskussion}
Die durchgeführten Experimente zeigen, dass der gewählte Gauß-basierte Ansatz grundsätzlich für die lokale Hindernisvermeidung geeignet ist. Im Vergleich zu globalen Planungsverfahren wie A* oder SLAM-basierten Navigationslösungen besitzt der entwickelte Ansatz deutlich geringere Anforderungen an Rechenleistung und Speicherbedarf. Dadurch eignet er sich insbesondere für ressourcenbeschränkte Systeme wie den Raspberry Pi. Gleichzeitig ergeben sich Einschränkungen. Da ausschließlich lokale Sensordaten betrachtet werden, besitzt das Fahrzeug keine Kenntnis über die globale Umgebung. Komplexe Hindernisanordnungen können daher nicht optimal berücksichtigt werden. Die in den Testläufen erreichte mittlere Erfolgsquote von 90\,\% bei den Ausweichmanövern zeigt, dass der entwickelte Ansatz grundsätzlich für autonome Fahrsysteme im Modellmaßstab geeignet ist. Die Hinderniserkennung erreichte über alle Testszenarien hinweg eine Erfolgsquote von 98\,\%. Für einen produktiven Einsatz wären jedoch zusätzliche Verfahren zur Kartierung, Lokalisierung und globalen Pfadplanung erforderlich.